\begin{blocksection}
\question Implement \lstinline{tree_height}, which takes in a tree and returns its height: the depth of the lowest leaf, where the root has depth 0.

\begin{lstlisting}
def tree_height(t: Tree[int]) -> int:
    """
    >>> tree_height(Tree(1))
    0
    >>> t = Tree(1, [Tree(4), Tree(2, [Tree(3)])])
    >>> tree_height(t)
    2
    """
    if ______________________________________:

        return ______________________________

    return 1 + ________________________________________________________
\end{lstlisting}

\begin{solution}[1in]
\begin{lstlisting}
def tree_height(t: Tree[int]) -> int:
    if is_leaf(t):
        return 0
    return 1 + max([tree_height(b) for b in t.branches])
\end{lstlisting}
\end{solution}

\end{blocksection}

\begin{questionmeta}
\textbf{Teaching Tips}\par\nopagebreak[4]
Suggested Time: 5 min; Difficulty: Easy
\nopagebreak[4]
\begin{itemize}
    \item \textbf{Use the lecture framework:} What small choice do we make (which branch?), what is the recursive call for each option, and how do we combine the results? Here, call \lstinline{tree_height} on every branch and take the \lstinline{max}.
    \item \textbf{Why the base case matters:} \lstinline{max} of an empty list errors, so leaves must be handled separately.
\end{itemize}
\end{questionmeta}
